\documentclass{article}
\usepackage{amsmath}
\usepackage{amssymb}
\usepackage{array}
\usepackage{booktabs}
\usepackage{braket}

\begin{document}

\section*{Monic LP: A Middle Ground for Pauli Reduction}

\subsection*{1. The Problem with the Two Extremes}

Experiments~1--3 identified two extreme strategies for choosing the collocation
matrix:

\begin{itemize}
    \item \textbf{Case A (monomial basis).}  The matrix $A_{\mathrm{mono}}$ saturates
          all $4^n$ Pauli strings.  For $N=4$ this means 16 strings and
          $L^2(n+1) = 768$ Hadamard tests per VQLS gradient step; for $N=8$ it means
          64 strings and $16{,}384$ tests per step.

    \item \textbf{Global LP ($M^* = I_N$).}  The unrestricted LP over all $N^2$
          entries of $M$ always finds $M^* = I_N$, requiring only 1 Pauli string and
          $n+1$ tests per step.  However, $M^* = I_N$ turns the VQLS system
          $I_N\ket{\psi} = \ket{b}$ into the trivial problem $\ket{\psi} = \ket{b}$,
          which is solvable classically (or by a single state-preparation circuit) and
          offers no quantum advantage.
\end{itemize}

This raises a natural question: \emph{can we reduce the Pauli count substantially,
while keeping the quantum system non-trivial?}  The monic LP provides an affirmative
answer.

\subsection*{2. The Monic LP Formulation}

Recall the basis-change identity $M = A_{\mathrm{mono}}\,C^{\top}$, where $C$ is the
coefficient matrix of the polynomial basis ($C_{jk}$ is the coefficient of $t^k$ in
$\phi_j$).  The monomial basis corresponds to $C = I_N$.

\paragraph{Restriction.}
We restrict $C$ to be \emph{lower-triangular with unit diagonal}:
\begin{equation}
    C_{jj} = 1 \text{ for all } j,
    \qquad
    C_{jk} = 0 \text{ for } k > j,
    \qquad
    C_{jk} \text{ free for } k < j.
\end{equation}
This is called a \emph{monic} basis: each $\phi_j$ is a degree-$j$ polynomial with
leading coefficient~1.  The free variables are the $\tfrac{N(N-1)}{2}$ strictly
lower-triangular entries.

\paragraph{Linearity.}
Writing $C = I + \Delta$ where $\Delta$ is strictly lower-triangular, the collocation
matrix becomes $M = A_{\mathrm{mono}}\,(I+\Delta)^{\top} = A_{\mathrm{mono}} + A_{\mathrm{mono}}\,\Delta^{\top}$.
Since each Pauli coefficient is $c_P = \tfrac{1}{N}\operatorname{Tr}(P\,M)$ and $M$
is affine in $\Delta$, we have
\begin{equation}
    c_P(\Delta) = c_P^0 + \mathbf{j}_P^{\top}\,\mathbf{x},
\end{equation}
where $c_P^0$ are the Case~A Pauli coefficients (the baseline), $\mathbf{x}$ is the
vector of free entries of $\Delta$, and $\mathbf{j}_P$ is a row of the Jacobian matrix
$J$ computed from the monomial collocation matrix.

\paragraph{LP.}
Introduce slack variables $t_P \geq |c_P|$ and minimise:
\begin{equation}
    \min_{\mathbf{x},\,\mathbf{t}}\;\sum_{P} t_P
    \qquad\text{s.t.}\quad
    J\mathbf{x} + \mathbf{c}^0 \leq \mathbf{t},
    \quad
    -J\mathbf{x} - \mathbf{c}^0 \leq \mathbf{t},
    \quad
    \mathbf{t} \geq 0,
    \quad
    \mathbf{x} \text{ free}.
    \label{eq:monic-lp}
\end{equation}
The LP has $\tfrac{N(N-1)}{2} + 4^n$ variables and $2\cdot 4^n$ inequality constraints.

\paragraph{No normalization constraint required.}
The global LP needs the equality $\operatorname{Tr}(M) = N$ to prevent $M \to 0$.
The monic LP does not: fixing $C_{jj}=1$ makes $C$ always invertible, so $M =
A_{\mathrm{mono}}\,C^{\top}$ is always invertible (trace cannot be made zero).
This is a structural simplification relative to the global LP.

\subsection*{3. Results}

Table~\ref{tab:monic} shows the Pauli counts and VQLS Hadamard-test overhead for
all three experiments.

\begin{table}[h]
\centering
\caption{Pauli strings and Hadamard tests per VQLS step ($\mathrm{HT/step} = L^2(n+1)$).}
\label{tab:monic}
\begin{tabular}{lccccc}
\toprule
Experiment & $N$ & $n$ & Case A & Monic LP & Global LP \\
\midrule
Exp.~1: $u''=0$ on $[0,1]$                  & 4 & 2 & 16/16 & \textbf{6/16} & 1/16 \\
Exp.~2: $u''+u'+u=f$ on $[-1,1]$            & 4 & 2 & 16/16 & \textbf{10/16} & 1/16 \\
Exp.~3: $u''+u'+u=f$ on $[-1,1]$ (scaled)  & 8 & 3 & 64/64 & \textbf{36/64} & 1/64 \\
\midrule
 & & & & & \\[-8pt]
Exp.~1 HT/step & & & 768 & \textbf{108} & 3 \\
Exp.~2 HT/step & & & 768 & \textbf{300} & 3 \\
Exp.~3 HT/step & & & 16{,}384 & \textbf{5{,}184} & 4 \\
\bottomrule
\end{tabular}
\end{table}

\paragraph{Monic basis for Experiment~1.}
The LP returns the basis coefficient matrix
\begin{equation}
C^* =
\begin{bmatrix}
1 & 0 & 0 & 0 \\
\tfrac{1}{2} & 1 & 0 & 0 \\
0 & -1 & 1 & 0 \\
0 & \tfrac{1}{4} & -\tfrac{5}{4} & 1
\end{bmatrix},
\end{equation}
giving six nonzero Pauli strings: $\{II, IX, IY, ZI, ZY, ZZ\}$.  The corresponding
collocation system is verified to recover the exact solution $u(t) = t$ at all test
points.

\paragraph{Why monic LP cannot reach the 3-Pauli Case~B solution.}
Experiment~1 showed that the analytical block-diagonal strategy achieves 3 Pauli
strings with a basis whose coefficient matrix has diagonal $(C_{00}, C_{11}, C_{22},
C_{33}) = (2, 2, 1, 1)$.  Normalising this to a unit diagonal by dividing each row
by its leading coefficient gives a monic matrix, but the resulting $M$ has
\emph{six} nonzero Pauli strings---identical to the monic LP optimum.  The 3-Pauli
solution is therefore inaccessible within the unit-diagonal family; it requires a
non-monic scaling of the diagonal.  This confirms that the monic LP searches a
strictly smaller space than the analytical Case~B construction.

\paragraph{Experiment~3 basis.}
For $N=8$ the monic LP has 28 free variables and finds 36 out of 64 Pauli strings,
reducing HT/step from 16{,}384 to 5{,}184.  All 36 nonzero Paulis are
three-qubit strings; the system is verified numerically at five test points.

\subsection*{4. Is There a Genuine Quantum Speedup?}

\paragraph{Hadamard-test reduction.}
The VQLS local cost function requires $L^2(n+1)$ Hadamard tests per gradient
evaluation, where $L$ is the number of nonzero Pauli strings.  Monic LP reduces $L$
from 16 to 6 for Experiment~1 (a $7.1\times$ reduction in HT/step) and from 64 to
36 for Experiment~3 (a $3.2\times$ reduction).  These are genuine reductions in
quantum circuit evaluations.

\paragraph{Non-triviality.}
Crucially, the monic LP result $M^* \neq I_N$, so the VQLS system $M^*\ket{\psi}
= \ket{b}$ is \emph{not} classically trivial.  Unlike the global LP, the monic
LP does not collapse the quantum problem into a state-preparation task.

\paragraph{Absolute vs.\ relative advantage.}
For the small sizes $N = 4, 8$ studied here, any $4\times 4$ or $8\times 8$ linear
system is trivially solvable classically in $O(1)$ time, so no VQLS formulation at
this scale offers a quantum advantage.  The question of quantum speedup only becomes
meaningful at much larger $N$, where quantum linear solvers are hypothesised to offer
an exponential advantage over classical methods (e.g.\ HHL with $O(\log N)$ depth).
At those scales the Hadamard-test reduction from the monic LP---a polynomial factor of
$L_A^2 / L_{\mathrm{monic}}^2$---would translate directly to a constant-factor
speedup in the quantum circuit overhead.

\paragraph{Convergence.}
VQLS was run for 150 gradient steps on each monic-LP system.  For Experiment~1 the
cost function converges to zero (the ansatz can represent $\ket{b}$).  For
Experiments~2 and~3 the cost plateaus near $0.025$ and $0.012$ respectively---the
same plateau observed for the global LP.  This plateau is not caused by Pauli
overhead; it is an \emph{ansatz limitation}: the single-layer $R_Y$ circuit with $n$
parameters cannot express the exact solution state $\ket{b}$ for general $b$.  Both
LP strategies share this bottleneck equally.

\subsection*{5. Is the Monic LP Easier to Solve?}

\paragraph{Problem size.}
\begin{center}
\begin{tabular}{lcc}
\toprule
 & Monic LP & Global LP \\
\midrule
Free variables    & $\tfrac{N(N-1)}{2}$ & $N^2$ \\
Slack variables   & $4^n = N^2$         & $4^n = N^2$ \\
Equality constraints & 0                & 1 ($\operatorname{Tr}(M)=N$) \\
Inequality constraints & $2\cdot 4^n$   & $2\cdot 4^n$ \\
\midrule
Total variables ($N=4$) & 22 & 32 \\
Total variables ($N=8$) & 92 & 128 \\
\bottomrule
\end{tabular}
\end{center}
The monic LP is modestly smaller---approximately $\tfrac{3}{4}$ the size of the
global LP in terms of variables---but the difference is not asymptotically significant.
Both are LPs of size $O(N^2)$ and both are solved in milliseconds by the HiGHS
simplex method at the scales considered here.

\paragraph{Structural simplicity.}
The monic LP is structurally simpler in one important respect: it requires
\emph{no normalization constraint}.  The unit-diagonal fixing ensures the basis is
always non-degenerate, so the scale of $M$ cannot collapse to zero.  The global LP
requires the explicit equality $\operatorname{Tr}(M) = N$ to prevent this collapse;
without it, the LP drives all Pauli coefficients to zero by shrinking $M \to 0$.
The monic formulation avoids this degeneracy by construction.

\paragraph{Solution quality.}
The global LP uniformly achieves 1 Pauli string regardless of the operator or domain.
The monic LP is operator-dependent and consistently worse:
\begin{center}
\begin{tabular}{lccc}
\toprule
 & Case A & Monic LP & Global LP \\
\midrule
Pauli strings ($N=4$, Exp.~1) & 16 & 6 & 1 \\
Pauli strings ($N=4$, Exp.~2) & 16 & 10 & 1 \\
Pauli strings ($N=8$, Exp.~3) & 64 & 36 & 1 \\
\bottomrule
\end{tabular}
\end{center}
The global LP's uniform optimum (always $M^* = I_N$, always 1 Pauli) is a consequence
of the fact that $I_N$ minimises $\sum_P|c_P|$ subject to $\operatorname{Tr}(M)=N$ for
\emph{any} matrix: non-identity Paulis are traceless, so $c_P(I_N)=0$ for all $P\neq
I^{\otimes n}$.  The monic LP, by contrast, cannot reach $I_N$ (since $C$ must be
lower-triangular with unit diagonal, and $A_{\mathrm{mono}}^{-T}$ is not of this form
in general), so its optimum depends on the specific operator and collocation points.

\subsection*{6. Summary}

The monic LP is a genuine middle ground between Case~A and the global LP.

\begin{itemize}
    \item \textbf{Easier to formulate} than the global LP: no normalization
          constraint, fewer variables, and a natural geometric interpretation (each
          basis function is a monic polynomial).
    \item \textbf{Not dramatically easier computationally}: both are $O(N^2)$-variable
          LPs solved in milliseconds.  The monic LP has $\tfrac{N(N-1)}{2}$ vs.\
          $N^2$ free variables, a factor of $\approx 2$ difference.
    \item \textbf{Provides moderate Hadamard-test speedup}: $7\times$ for Experiment~1,
          $2.6\times$ for Experiment~2, $3.2\times$ for Experiment~3.
    \item \textbf{Preserves quantum non-triviality}: $M^* \neq I_N$, so the quantum
          task is not classically solvable by inspection.
    \item \textbf{Cannot recover all analytical solutions}: the 3-Pauli Case~B
          solution from Experiment~1 requires non-unit diagonal coefficients and lies
          outside the monic family.  The monic LP optimum for Experiment~1 is 6
          Pauli strings, not 3.
\end{itemize}

The global LP remains the right choice if minimising Pauli count is the sole objective,
but its result ($M^* = I_N$) is classically trivial and therefore not useful for
demonstrating quantum advantage.  The monic LP offers a principled way to reduce
quantum circuit overhead while preserving a non-trivial quantum computation.  Whether
this overhead reduction is sufficient to yield a practical quantum advantage depends
on the problem size and the available ansatz depth---questions left for future work.

\end{document}